% !TEX root = ../main.tex
\section{The tableau}

We follow Section~2 of \cite{rupprecht2026sparse}. A sparse target state is
\[
  \ket{\Theta} = \sum_{i=0}^{s-1} c_i \ket{C_i},
\]
with each $\ket{C_i}$ an $n$-bit computational basis state and $s \ll 2^n$.

\begin{definition}[Tableau]
Stack the bitstrings $C_i$ as the rows of an $s \times n$ bit matrix.
Let $l = \ceillog{s}$. Split every row at position $l$ into
$\ket{C_i} = \ket{C'_i}_{s}\ket{C''_i}_{r}$, the \emph{subspace register}
($l$ qubits) and the \emph{non-subspace register} ($n-l$ qubits).
\end{definition}

The goal is a circuit $G$ and an injective map $f$ with
$G\ket{C_i} = \ket{f(i)}_s \ket{0}_r$ for all $i$. The state preparation
circuit is then the inverse of $G$ applied after a dense preparation on
$l$ qubits.

\begin{definition}[Batch size]
$m = \max\{\,2^p : 2^p \le n - l\,\}$.
\end{definition}

\paragraph{Worked example.} The paper's seven states have $s=7$, $n=7$,
so $l=\ceillog{7}=3$ and $m=4$. The first state is
\bits{011|1011}: subspace part \bits{011}, non-subspace part \bits{1011}.

\begin{remark}
Every gate in $G$ (X, CX, Toffoli, multi-controlled X) permutes
computational basis states, so $G$ can be simulated on bitstrings alone.
This is why the code can be tested classically.
\end{remark}